For psh functions \(u_1,\ldots,u_p\) on a coordinate open set \(U\), put
\[
 O_k=\bigcap_{j=1}^p\{u_j>-k\}.
\]
On \(O_k\), Bedford--Taylor theory and plurifine locality give
\[
 \one_{O_k}\ddc u_1\wedge\cdots\wedge\ddc u_p
 =
 \one_{O_k}\bigwedge_{j=1}^p\ddc\max\{u_j,-k\}.
\tag{51.15}\label{eq:complex-npp-truncation}
\]
The positive Borel \((p,p)\)-valued measures on the right increase with
\(k\).  Their increasing limit need not yet be a current: its coefficients
can have infinite mass on every neighbourhood of a point.  Fix a smooth
strictly positive \((1,1)\)-form \(\omega_U\) on \(U\).  We say that the
local nonpluripolar product is \emph{well defined} when, for every compact
set \(K\Subset U\),
\[
 \sup_k\int_{K\cap O_k}\omega_U^{n-p}\wedge
       \bigwedge_{j=1}^p\ddc\max\{u_j,-k\}<+\infty.
\]
This condition does not depend on \(\omega_U\), since any two strictly
positive forms are uniformly comparable on \(K\).  Under this condition the
increasing limit has locally finite coefficients and hence defines a positive
\((p,p)\)-current, denoted by
\[
              \left\langle\ddc u_1\wedge\cdots\wedge\ddc u_p\right\rangle .
\]
Without the condition, the notation denotes at most a locally infinite
positive Borel object and is not used as a current in this book.  The
qualification is essential for \(p\geq2\): Kiselman's example, recalled in
\cite[Example~1.3]{BEGZ10}, gives a psh function on a domain in
\(\mathbf C^2\) for which the measure obtained from the smooth products off
its polar axis is not locally finite.

For global closed positive currents
\(T_j=\theta_j+\ddc\varphi_j\), choose local psh potentials of \(T_j\).
The comparison below proves that local finiteness is independent of those
choices and that the resulting local currents agree on overlaps.  They
therefore patch to a global current, denoted by
\[
       \langle T_1\wedge\cdots\wedge T_p\rangle .
\tag{51.16}\label{eq:complex-global-npp}
\]
The compact K\"ahler case is special: the local mass condition is automatic.
This is proved, rather than assumed, in
\Cref{thm:complex-npp-ledger} below.

We record the three details of the definition that will be used in that
proof.

\emph{Compatibility.}  On the plurifine open set \(O_k\), each \(u_j\)
coincides with \(\max\{u_j,-k\}\).  Plurifine locality of the
Bedford--Taylor product therefore gives
\eqref{eq:complex-npp-truncation}.  Applying the same observation with two
levels shows that the restricted products increase with \(k\).

\emph{Admissible exhaustions.}  Suppose that the product \(R\) is locally
finite.  Let \(a_k\) be an increasing sequence of positive integers tending
to infinity, put
\[
 O_{a_k}=\bigcap_{j=1}^p\{u_j>-a_k\},
\]
and let \(E_k\) be increasing Borel sets such that
\(E_k\subset O_{a_k}\) and
\(U\setminus\bigcup_kE_k\) is pluripolar.  The current at the \(k\)-th
stage is
\[
 R_k^{E}:=\one_{E_k}\bigwedge_{j=1}^p
                    \ddc\max\{u_j,-a_k\}.
\tag{51.15a}\label{eq:complex-npp-admissible-exhaustion}
\]
Plurifine locality gives \(R_k^{E}=\one_{E_k}R\).  The no-charge assertion
proved in \Cref{thm:complex-npp-ledger} therefore implies
\(R_k^{E}\to R\), weakly and against bounded measurable coefficient test
forms with compact support.  Equivalently, this is dominated convergence for
the locally finite trace measure of \(R\).  This precise form of
\cite[Lemma~1.5]{BEGZ10} is the one used below.  The capacity proof of no
pluripolar charge precedes every use of this rule in the proof below, so there
is no circularity.

\emph{Dependence only on the currents.}  Let \(u'_j\) be another psh
potential of the same current on an open set \(W\), and write
\(h_j=u'_j-u_j\); thus \(h_j\) is pluriharmonic.  Fix
\(K\Subset V\Subset W\), and choose a positive integer \(C\) such that
\[
                         |h_j|\leq C
\]
on a neighbourhood of \(\overline V\), for every \(j\).  Set
\(O'_k=\bigcap_j\{u'_j>-k\}\).  The inclusion
\[
 O'_k\cap V\subset O_{k+C}\cap V
       =\bigcap_{j=1}^p\{u_j>-k-C\}\cap V
\tag{51.15b}\label{eq:complex-npp-potential-locus-forward}
\]
and plurifine locality give
\[
 \one_{O'_k\cap V}\bigwedge_{j=1}^p\ddc\max\{u'_j,-k\}
 =\one_{O'_k\cap V}\bigwedge_{j=1}^p\ddc\max\{u_j,-k-C\}.
\tag{51.15c}\label{eq:complex-npp-potential-current-forward}
\]
Indeed, on \(O'_k\cap V\) the functions on the left equal \(u_j+h_j\),
whereas \(u_j>-k-C\); hence both sides there are the Bedford--Taylor product
of the currents \(\ddc u_j=\ddc(u_j+h_j)\).  Consequently,
\[
\begin{aligned}
 &\int_{K\cap O'_k}\omega_U^{n-p}\wedge
       \bigwedge_{j=1}^p\ddc\max\{u'_j,-k\}\\
 &\qquad\leq
 \int_{K\cap O_{k+C}}\omega_U^{n-p}\wedge
       \bigwedge_{j=1}^p\ddc\max\{u_j,-k-C\}.
\end{aligned}
\tag{51.15d}\label{eq:complex-npp-potential-mass-forward}
\]
The right-hand side is bounded by the defining supremum for the unprimed
tuple.  Thus local finiteness for the unprimed tuple implies it for the
primed tuple.  Using \(-h_j\) gives the reverse inclusion
\(O_k\cap V\subset O'_{k+C}\cap V\) and the reverse comparison
\[
 \int_{K\cap O_k}\omega_U^{n-p}\wedge
       \bigwedge_{j=1}^p\ddc\max\{u_j,-k\}
 \leq
 \int_{K\cap O'_{k+C}}\omega_U^{n-p}\wedge
       \bigwedge_{j=1}^p\ddc\max\{u'_j,-k-C\}.
\tag{51.15e}\label{eq:complex-npp-potential-mass-reverse}
\]
The right-hand side is now bounded by the defining supremum for the primed
tuple.  This proves both directions of potential-choice invariance for the
local mass condition.

Assume now that this equivalent condition holds.  On \(V\) define the
increasing common exhaustion
\[
                         E_k=O_k\cap O'_k\cap V.
\]
The set \(V\setminus\bigcup_kE_k\) is contained in the finite union of the
pole sets of the \(u_j\) and \(u'_j\), hence is pluripolar.  On \(E_k\),
plurifine locality and \(\ddc h_j=0\) give the termwise identity
\[
 \one_{E_k}\bigwedge_{j=1}^p\ddc\max\{u_j,-k\}
 =\one_{E_k}\bigwedge_{j=1}^p\ddc\max\{u'_j,-k\}.
\tag{51.15f}\label{eq:complex-npp-potential-common-current}
\]
After the capacity argument in \Cref{thm:complex-npp-ledger} has established
that the two locally finite limits charge no pluripolar set,
\eqref{eq:complex-npp-admissible-exhaustion} applied to both sides of
\eqref{eq:complex-npp-potential-common-current} identifies those limits on
\(V\).  Since \(K\Subset V\Subset W\) was arbitrary, the local products patch
and depend only on the currents.  This proves the second assertion of
\cite[Proposition~1.4]{BEGZ10} with the mass comparisons included.
